\documentclass[11pt]{article}

\usepackage[T1]{fontenc}
\usepackage[utf8]{inputenc}
\usepackage[english,turkish]{babel}
\usepackage{amsmath,amssymb}
\usepackage{booktabs}
\usepackage{array}
\usepackage{geometry}
\usepackage{longtable}
\usepackage{hyperref}
\usepackage{algorithm}
\usepackage{algpseudocode}

\geometry{margin=1in}

\title{MWCCVC-Based Energy-Efficient Secure Link Monitoring in Wireless Sensor Networks for Industrial IoT\\Bilingual Draft (English--Turkish)}
\author{}
\date{}

\begin{document}

\maketitle
\tableofcontents
\clearpage

\selectlanguage{english}

\section*{Step 1: Implementations Found}

The five active implementations identified in the repository are as follows.

\begin{itemize}
    \item Greedy-1: \texttt{GCCVC}, entry point \texttt{solve\_gccvc}, file \texttt{backend/app/algorithms.py}.
    \item Greedy-2: \texttt{GRCCVC}, entry point \texttt{solve\_grccvc}, file \texttt{backend/app/algorithms.py}.
    \item Greedy-3: \texttt{GWCCVC}, entry point \texttt{solve\_gwccvc}, file \texttt{backend/app/algorithms.py}.
    \item HGA-1: \texttt{HGA-CCVC}, entry point \texttt{solve\_hga}, file \texttt{backend/app/algorithms.py}.
    \item HGA-2: \texttt{HGA-CCVC v2}, class \texttt{HGA\_CCVC\_V2} and wrapper \texttt{solve\_hga\_v2}, file \texttt{backend/app/algorithms.py}.
\end{itemize}

\section*{Step 2: Dataset Split and Experimental Protocol Extracted from Code}

The repository uses a single dataset source located in \texttt{backend/data/DagdevirenDataset}. Graph files follow the pattern \texttt{n\{n\}\_m\{m\}\_s\{s\}.txt}. The default scale split defined in code is
\[
\text{Small}=\{10,15,20,25\},\qquad
\text{Medium}=\{50,100,150,200\},\qquad
\text{Large}=\{250,500,750,1000\}.
\]
The default connectivity-ratio filter is
\[
\frac{m}{n}\in\{2,4,6,8\}.
\]
The default scale-wise capacity values are
\[
K_{\text{small}}=18,\qquad K_{\text{medium}}=16,\qquad K_{\text{large}}=16.
\]
The preset selection policy includes all available seed indices \(s\) for every selected \((n,m)\) pair. According to the corrected experimental description used for the manuscript, the testing pipeline was organized in two stages. First, for each graph instance, the monitor capacity \(K\) was selected automatically. The search interval was bounded by
\[
K_{\min}=\left\lceil \frac{|E|}{|V|}\right\rceil,\qquad K_{\max}=\Delta(G).
\]
Within this interval, a binary search was used to identify the minimum feasible \(K\), and the feasibility/weight selector for this stage was the ratio-based greedy method \texttt{GRCCVC}. When the optimization target was minimum weight, the procedure then performed a sampled scan over feasible \(K\) values and selected the \(K\) that produced the minimum valid \texttt{GRCCVC} weight, with ties broken in favor of smaller \(K\). Second, after fixing this instance-specific \(K\), the five methods were run on the same graph using that selected capacity. In the reported configuration, the population size and generation count were set to \texttt{100} and \texttt{100}, respectively.

\paragraph{Integrity note.}
No results are inserted in this file. Small-scale and medium-scale result tables are intentionally left as placeholders to be filled from the external DOCX. Large-scale evaluation is explicitly marked as planned/pending.

\setcounter{section}{2}

\section{System Model}

The industrial wireless sensor network (WSN) is modeled as an undirected graph \(G=(V,E)\), where each vertex \(v\in V\) denotes a sensor or monitor-capable node and each edge \(e=(u,v)\in E\) denotes a communication link. In the implementation, nodes carry an identifier, a weight, and a monitoring capacity, whereas edges carry a covered/uncovered status and a record of the node to which the edge has been assigned for monitoring. Thus, monitoring is endpoint-based.

A link \(e=(u,v)\) is considered monitored if at least one of its endpoints is selected into the monitoring set and the edge is assigned to that selected endpoint. Let \(C\subseteq V\) denote the selected monitor set. Then every \(e=(u,v)\in E\) must satisfy \(u\in C\) or \(v\in C\), and each covered edge is associated with one selected incident endpoint.

The node-weight term \(w(v)\) is taken directly from the dataset as an exogenous scalar cost. The code does not define an explicit transformation from residual battery, radio duty cycle, or security exposure to this weight. Therefore, the most faithful interpretation is that \(w(v)\) is a generic activation cost used to model monitoring burden or energy expenditure. The optimization objective minimizes the total weight of the selected connected monitoring set.

Each node receives the same per-experiment capacity value \(K\), i.e.,
\[
\mathrm{cap}(v)=K,\qquad \forall v\in V.
\]
This capacity limits the number of incident links assigned to node \(v\):
\[
\sum_{e\in \delta(v)} y_{e,v} \le K.
\]
The implementation maintains a residual capacity value for each node and decrements it whenever the node claims an incident edge.

The connectivity requirement reflects a connected monitoring backbone. Feasibility therefore requires that the subgraph induced by \(C\) be connected. In the present repository, secure monitoring is represented at the graph level as complete link observability by a connected, capacity-feasible set of monitoring nodes. Packet inspection, IDS signatures, cryptographic protocols, adversarial models, and traffic semantics are not explicitly modeled in the provided artifacts.

\section{Problem Formulation (MWCCVC)}

Let \(x_v\in\{0,1\}\) indicate whether node \(v\in V\) is selected as a monitor. To model exact edge-to-monitor assignment under node capacity, let \(y_{e,v}\in\{0,1\}\) be defined only when node \(v\) is an endpoint of edge \(e\). The optimization problem solved by the repository can be written as
\[
\min \sum_{v\in V} w(v)x_v.
\]

The edge-coverage constraint is
\[
x_u+x_v\ge 1,\qquad \forall (u,v)\in E.
\]
Under the assignment view, each edge must be assigned to exactly one selected endpoint:
\[
\sum_{v\in e} y_{e,v}=1,\qquad \forall e\in E,
\]
with endpoint consistency
\[
y_{e,v}\le x_v,\qquad \forall e\in E,\; \forall v\in e.
\]

The capacity constraint is
\[
\sum_{e\in \delta(v)} y_{e,v}\le Kx_v,\qquad \forall v\in V.
\]
This formulation matches the implementation closely because capacity feasibility is checked through an exact endpoint-assignment routine based on forced assignments and a max-flow validation for flexible edges.

Connectivity is imposed by requiring the induced subgraph \(G[C]\), where \(C=\{v\in V:x_v=1\}\), to be connected:
\[
G[C]\text{ is connected.}
\]
The code does not enforce connectivity via mixed-integer cuts. Instead, it constructs a candidate cover and then repairs disconnected covers by inserting bridge vertices along shortest connecting paths.

The resulting problem is a weighted capacitated connected vertex cover variant. It is NP-hard because classical vertex cover and connected vertex cover appear as special cases [CITATION NEEDED]. Consequently, the repository relies on greedy heuristics and hybrid genetic algorithms.

\section{Proposed Methods}

All five methods target the same feasibility conditions: full edge coverage, connectedness of the selected set, and exact capacity feasibility. The three greedy methods differ in their node-priority and edge-ordering rules. The two HGA variants search over parameterized decoders rather than directly encoding a final cover.

The implementation also uses a color-based guidance mechanism: \texttt{WHITE} for untouched nodes, \texttt{BLACK} for selected monitors, \texttt{GRAY} for nodes adjacent to selected monitors, and \texttt{RED} for priority nodes identified by the marking rule. These colors are heuristic search states rather than decision variables of the formal optimization model.

\subsection{Greedy-1}

Greedy-1 corresponds to \texttt{GCCVC}. It is degree-oriented: among eligible candidates, it selects the node with maximum uncovered degree, breaking ties in favor of smaller node weight. Once a node is selected, its uncovered incident edges are ordered by decreasing uncovered degree of the neighboring endpoint, then by increasing neighbor weight. This makes the algorithm strongly coverage-oriented both in node selection and in local edge assignment.

Coverage is enforced incrementally by assigning uncovered edges to selected nodes while residual capacity remains positive. If uncovered edges still exist after the main constructive phase, the method uses a fallback completion stage that first consumes the remaining capacities of already selected nodes and then selects additional nodes as needed.

Connectivity is enforced by a capacity-aware repair. If the selected set induces multiple components, the method computes low-cost bridge paths and joins components through a union-find-guided bridging process. The repair adds path vertices to the cover; it does not remove already selected vertices.

\begin{algorithm}[h]
\caption{Greedy-1 / GCCVC}
\begin{algorithmic}[1]
\State \textbf{Input:} Graph \(G=(V,E)\), weights \(w(v)\), capacity \(K\)
\State \textbf{Output:} Connected capacitated cover \(C\)
\State \(C\gets \emptyset\); initialize all nodes as \texttt{WHITE}; set residual capacities to \(K\)
\While{there exists an uncovered edge}
    \State mark priority (\texttt{RED}) nodes using uncovered-degree rules
    \State build candidate pool from non-\texttt{BLACK} nodes with positive uncovered degree
    \State choose \(v\) maximizing uncovered degree; break ties by smaller weight
    \State add \(v\) to \(C\) and mark it \texttt{BLACK}
    \State sort uncovered incident edges of \(v\) by neighbor uncovered degree (descending), then neighbor weight
    \For{each sorted uncovered incident edge of \(v\)}
        \If{residual capacity of \(v\) is zero}
            \State \textbf{break}
        \EndIf
        \State assign the edge to \(v\) and decrease its residual capacity
    \EndFor
\EndWhile
\If{uncovered edges remain}
    \State force-cover the remaining edges
\EndIf
\State repair disconnected cover components by capacity-aware bridge paths
\State \Return \(C\)
\end{algorithmic}
\end{algorithm}

\subsection{Greedy-2}

Greedy-2 corresponds to \texttt{GRCCVC}. Its node-selection rule is ratio-based:
\[
\rho(v)=\frac{w(v)}{\sum_{u\in N_U(v)} w(u)},
\]
where \(N_U(v)\) is the currently uncovered neighborhood of \(v\). The method selects the node with minimum ratio, breaking ties in favor of lower node weight. Intuitively, it seeks nodes whose own activation cost is small relative to the aggregate cost structure of the uncovered neighborhood they can dominate.

Incident uncovered edges of the selected node are shuffled randomly before assignment. Hence, Greedy-2 combines a deterministic ratio-based node rule with randomized local edge ordering. Coverage, fallback completion, and connectivity repair follow the same general mechanism as in Greedy-1, including capacity-aware bridging.

\begin{algorithm}[h]
\caption{Greedy-2 / GRCCVC}
\begin{algorithmic}[1]
\State \textbf{Input:} Graph \(G=(V,E)\), weights \(w(v)\), capacity \(K\)
\State \textbf{Output:} Connected capacitated cover \(C\)
\State initialize \(C\gets\emptyset\) and residual capacities to \(K\)
\While{there exists an uncovered edge}
    \State mark \texttt{RED} nodes
    \State form a candidate pool of feasible non-\texttt{BLACK} nodes
    \For{each candidate \(v\)}
        \State compute \(\rho(v)=w(v)/\sum_{u\in N_U(v)}w(u)\)
    \EndFor
    \State choose \(v\) minimizing \(\rho(v)\); break ties by smaller weight
    \State add \(v\) to \(C\)
    \State randomly shuffle uncovered incident edges of \(v\)
    \State assign shuffled uncovered incident edges to \(v\) until its residual capacity is exhausted
\EndWhile
\If{uncovered edges remain}
    \State force-cover the remaining edges
\EndIf
\State repair disconnected cover components by capacity-aware bridge paths
\State \Return \(C\)
\end{algorithmic}
\end{algorithm}

\subsection{Greedy-3}

Greedy-3 corresponds to \texttt{GWCCVC}. It is weight-oriented: the method chooses the minimum-weight candidate, breaking ties by larger uncovered degree. It therefore emphasizes low activation cost more strongly than the other greedy variants.

After node selection, the local edge-ordering rule prioritizes uncovered edges incident to \texttt{WHITE} neighbors over those incident to non-\texttt{WHITE} neighbors. This favors expansion into previously untouched regions of the graph. Capacity is enforced directly during edge assignment.

Coverage completion again uses the common fallback mechanism. The connectivity repair, however, differs from the previous greedy variants. Greedy-3 uses a simpler unconstrained shortest-path repair: it selects a representative node from each component, specifically the minimum-weight node in that component, and then connects these representatives to a base component through BFS shortest paths.

\begin{algorithm}[h]
\caption{Greedy-3 / GWCCVC}
\begin{algorithmic}[1]
\State \textbf{Input:} Graph \(G=(V,E)\), weights \(w(v)\), capacity \(K\)
\State \textbf{Output:} Connected capacitated cover \(C\)
\State initialize \(C\gets\emptyset\) and residual capacities to \(K\)
\While{there exists an uncovered edge}
    \State mark \texttt{RED} nodes
    \State build the candidate pool
    \State choose \(v\) with minimum weight; break ties by maximum uncovered degree
    \State add \(v\) to \(C\)
    \State sort uncovered incident edges so that \texttt{WHITE} neighbors are prioritized
    \State assign sorted uncovered incident edges to \(v\) until its residual capacity is exhausted
\EndWhile
\If{uncovered edges remain}
    \State force-cover the remaining edges
\EndIf
\State identify connected components of \(G[C]\)
\State choose the minimum-weight representative of each component
\State connect representatives to a base component through BFS shortest paths
\State add bridge-path vertices to \(C\)
\State \Return \(C\)
\end{algorithmic}
\end{algorithm}

\subsection{HGA-1}

HGA-1 corresponds to \texttt{HGA-CCVC}. Its chromosome has two parts: a six-dimensional real-valued parameter vector
\[
\boldsymbol{\theta}=(a,b,c,d,e,\epsilon)\in[0,3]^6,
\]
and a node-wise random-key map. The chromosome does not directly encode the final cover. Instead, it parameterizes a greedy decoder that constructs a cover from scratch.

For a candidate node \(v\), the decoder score can be written as
\[
S(v)=a\deg_U(v)-bw(v)-c\rho(v)+d\mathbf{1}_{\mathrm{RED}}(v)+e\mathbf{1}_{\mathrm{GRAY}}(v)+0.3a\frac{\deg_U(v)}{\max(w(v),10^{-9})}-\epsilon\,\mathrm{key}(v),
\]
where \(\deg_U(v)\) is uncovered degree and \(\rho(v)\) is the same weight-to-neighborhood ratio used in the code. Thus, HGA-1 searches for a good balance among coverage gain, node cost, neighborhood structure, color bonuses, and random-key diversification.

Initialization is hybrid. The initial population is seeded using covers produced by the three greedy methods together with predefined \(\theta\)-templates. Additional individuals are generated by perturbing seeded templates or by pure random initialization.

Constraint handling is hybrid as well. The decoder first constructs and repairs a cover; then the fitness function applies explicit penalties when the candidate is not a valid capacitated connected cover. A concise representation of the penalized cost is
\[
f(C)=W(C)+10^6U(C)+10^5(1-I_{\mathrm{conn}}(C))+10^5(1-I_{\mathrm{cap}}(C)),
\]
where \(W(C)\) is total weight, \(U(C)\) is the number of uncovered edges, and \(I_{\mathrm{conn}}\), \(I_{\mathrm{cap}}\) are connectedness and capacity-feasibility indicators.

HGA-1 further includes pruning and weight-focused local search, elitism, tournament selection, crossover over \(\theta\), random-key recombination, adaptive mutation, and immigrant injection. Under stagnation, a fresh \texttt{GRCCVC}-based seed is reintroduced into the population.

\begin{algorithm}[h]
\caption{HGA-1 / HGA-CCVC}
\begin{algorithmic}[1]
\State \textbf{Input:} Graph \(G=(V,E)\), weights, capacity \(K\), population size \(P\), generations \(G_{\max}\)
\State \textbf{Output:} Best feasible connected capacitated cover \(C^{\star}\)
\State initialize a hybrid population of chromosomes \((\boldsymbol{\theta},\mathrm{key})\)
\State \(C^{\star}\gets\emptyset\), \(f^{\star}\gets +\infty\)
\For{generation \(=1\) to \(G_{\max}\)}
    \For{each chromosome in the population}
        \State decode the chromosome into a cover \(C\) using a parameterized greedy construction
        \State repair connectivity of \(C\)
        \State optionally apply local search and weight-focused local search
        \State evaluate penalized fitness \(f(C)\)
        \If{\(f(C)<f^{\star}\)}
            \State update \(C^{\star}\) and \(f^{\star}\)
        \EndIf
    \EndFor
    \State sort the population by fitness and keep elites
    \If{stagnation persists}
        \State inject a new greedy-seeded individual
    \EndIf
    \While{the new population is not full}
        \State select parents by tournament selection
        \State apply crossover to \(\boldsymbol{\theta}\) and random keys
        \State mutate \(\boldsymbol{\theta}\) and random keys
        \State add the offspring to the new population
    \EndWhile
\EndFor
\State \Return \(C^{\star}\)
\end{algorithmic}
\end{algorithm}

\subsection{HGA-2}

HGA-2 corresponds to \texttt{HGA-CCVC v2}. It keeps the same basic chromosome concept as HGA-1 but makes the evolutionary design more modular and parameterized. The default parameter set in code includes \texttt{elite\_size=2}, \texttt{tournament\_k=3}, \texttt{theta\_mut\_prob=0.2}, \texttt{theta\_mut\_sigma=0.2}, \texttt{key\_mut\_prob=0.02}, \texttt{local\_search\_prob=0.2}, \texttt{local\_search\_moves=7}, \texttt{local\_search\_trials=36}, \texttt{top\_l=3}, \texttt{repair\_mode=tier2}, \texttt{stall\_generations=10}, and \texttt{seed\_bias=0.3}.

Population initialization is hybrid and seeded. Four strategy templates are used: one based on \texttt{GCCVC}, one based on \texttt{GRCCVC}, and two based on \texttt{GWCCVC}. These seeded chromosomes are then augmented with random chromosomes until the target population size is reached.

The decoder remains score-driven, but HGA-2 introduces a top-\(L\) stochastic selection mechanism. Instead of always taking the single best scored candidate, it samples from the best few nodes under the default stochastic mode. A similar perturbation is applied to local edge ordering, increasing diversification while preserving score guidance.

Its penalized objective can be summarized as
\[
f(C)=W(C)+0.001|C|+\lambda_1U(C)+\lambda_2(1-I_{\mathrm{conn}}(C))+\lambda_3\phi_{\mathrm{cap}}(C),
\]
with \(\lambda_1=10^6\), \(\lambda_2=10^5\), and \(\lambda_3=10^5\). The extra term \(0.001|C|\) acts as a tie-breaker favoring smaller covers when weights are nearly equal.

Selection uses \(k\)-tournament selection, crossover is arithmetic on the \(\theta\)-vector and uniform on the random keys, and mutation combines Gaussian perturbation of \(\theta\) with random key resets. A feasible-only local search then attempts profitable node removals and light-node swaps.

\begin{algorithm}[h]
\caption{HGA-2 / HGA-CCVC v2}
\begin{algorithmic}[1]
\State \textbf{Input:} Graph \(G=(V,E)\), weights, capacity \(K\), population size \(P\), generations \(G_{\max}\)
\State \textbf{Output:} Best feasible connected capacitated cover \(C^{\star}\)
\State initialize the population from greedy-based seeds and random chromosomes
\For{generation \(=1\) to \(G_{\max}\)}
    \For{each chromosome \((\boldsymbol{\theta},\mathbf{r})\)}
        \State decode it using a score-guided greedy construction
        \State under stochastic mode, sample from the top-\(L\) candidates
        \State repair connectivity according to the configured repair mode
        \State optionally apply feasible-only local search pruning
        \State compute the penalized fitness
    \EndFor
    \State rank chromosomes by fitness
    \State preserve elites
    \If{stall threshold is reached}
        \State terminate the run
    \EndIf
    \While{the population is not full}
        \State select parents by \(k\)-tournament
        \State apply arithmetic crossover to \(\boldsymbol{\theta}\)
        \State apply uniform crossover to node keys
        \State mutate \(\boldsymbol{\theta}\) and keys
        \State insert the offspring
    \EndWhile
\EndFor
\State \Return the best decoded feasible cover
\end{algorithmic}
\end{algorithm}

\section{Algorithmic Complexity Analysis}

Let \(n=|V|\), \(m=|E|\), \(P\) be the population size, and \(G\) be the number of generations. The exact capacity-feasibility routine used by the code builds a flow network over flexible edges and selected cover nodes. Denoting this validation cost by \(T_{\mathrm{flow}}(n,m)\), a conservative upper bound is cubic in the flow-network size.

For all three greedy methods, the dominant constructive cost arises from repeated rescoring of candidate nodes, repeated incident-edge scans, and subsequent connectivity repair. A conservative bound for the constructive phase is \(O(nm)\).

For Greedy-1 and Greedy-2, the connectivity repair is capacity-aware and may evaluate multiple inter-component shortest bridge paths. A conservative bound is
\[
T_{\text{Greedy-1}}=O\bigl(nm+\kappa^2m\log n\bigr),
\]
\[
T_{\text{Greedy-2}}=O\bigl(nm+\kappa^2m\log n\bigr),
\]
where \(\kappa\) is the number of components before repair.

For Greedy-3, the BFS-based representative-connection repair is simpler, yielding the conservative bound
\[
T_{\text{Greedy-3}}=O\bigl(nm+\kappa(n+m)\bigr).
\]

For HGA-1, each chromosome evaluation includes greedy decoding, connectivity repair, local search, and capacity validation for feasible candidates. With abstract local-search effort \(L_1\), a conservative bound is
\[
T_{\text{HGA-1}}=O\Bigl(PG\bigl(nm+\kappa^2m\log n+L_1T_{\mathrm{flow}}(n,m)\bigr)\Bigr).
\]

For HGA-2, the same structure applies, but local search is more explicitly bounded by \texttt{local\_search\_moves} and \texttt{local\_search\_trials}. With bounded local-search effort \(L_2\), a conservative bound is
\[
T_{\text{HGA-2}}=O\Bigl(PG\bigl(nm+\kappa^2m\log n+L_2T_{\mathrm{flow}}(n,m)\bigr)\Bigr).
\]
These bounds describe the current implementation rather than an asymptotically optimized data-structure design.

\section{Dataset Description (Small/Medium/Large split)}

The study uses a single topology source stored in \texttt{backend/data/DagdevirenDataset}. Files are identified by \((n,m,s)\), where \(n\) is node count, \(m\) is edge count, and \(s\) is the seed/index encoded in the filename. The reader supports an optional header of the form \texttt{n m}; otherwise the counts are inferred from the filename.

After the optional header, the file body is structured as follows:
\begin{enumerate}
    \item \(n\) vertex records of the form \texttt{id weight};
    \item \(m\) edge records of the form \texttt{u v}.
\end{enumerate}

The code-level default scale split is
\[
\text{Small}=\{10,15,20,25\},\qquad
\text{Medium}=\{50,100,150,200\},\qquad
\text{Large}=\{250,500,750,1000\}.
\]
The default connectivity-ratio filter is
\[
\frac{m}{n}\in\{2,4,6,8\}.
\]
The preset selection policy includes all available seed files for every retained \((n,m)\) pair.

\begin{table}[h]
\centering
\caption{Dataset Format}
\label{tab:dataset-format}
\begin{tabular}{lll}
\toprule
Field & Meaning & Example \\
\midrule
\texttt{n100\_m200\_s1.txt} & Filename encodes \((n,m,s)\) & \texttt{n100\_m200\_s1.txt} \\
\texttt{n m} & Optional header with declared counts & \texttt{100 200} \\
\texttt{id weight} & Vertex record & \texttt{0 116.20978281782678} \\
\texttt{u v} & Edge record & \texttt{0 1} \\
\bottomrule
\end{tabular}
\end{table}

\section{Experimental Setup and Protocol}

All methods are executed on the same graph instances within a scale bucket, with identical node weights, the same edge set, and the same capacity value selected for that instance. Feasibility is checked uniformly through the same verification routine, which reports edge coverage, connectedness, total selected weight, cover size, and exact endpoint-assignment capacity feasibility.

The five-method evaluation pipeline is
\[
\{\texttt{gccvc},\texttt{grccvc},\texttt{gwccvc},\texttt{hga},\texttt{hga\_v2}\}.
\]
In the corrected protocol, \(K\) is not fixed a priori by scale. Instead, for each graph instance, the search interval \([K_{\min},K_{\max}]\) is formed using
\[
K_{\min}=\left\lceil \frac{|E|}{|V|}\right\rceil,\qquad K_{\max}=\Delta(G).
\]
The repository then uses only \texttt{GRCCVC} in the \(K\)-selection stage. First, a binary search identifies the minimum feasible \(K\). Next, under the minimum-weight setting, a sampled scan is carried out over feasible \(K\) values between the minimum feasible \(K\) and \(K_{\max}\). The selected capacity is the one that yields the minimum valid \texttt{GRCCVC} weight, with ties broken by smaller \(K\). After this instance-specific \(K\) is fixed, all five methods are executed on the same graph using that \(K\). In the reported experiments, the configured values were \texttt{popSize=100}, \texttt{generations=100}, and a common seed setting.

\subsection{Small-Scale Setup}

The small-scale bucket contains \(n\in\{10,15,20,25\}\) and ratio filter \(m/n\in\{2,4,6,8\}\). All available seed files are selected for each retained \((n,m)\) pair. For each selected small-scale instance, \(K\) is first optimized using the \texttt{GRCCVC}-guided binary-search-plus-weight-scan procedure described above. After the instance-specific \(K\) is fixed, Greedy-1, Greedy-2, Greedy-3, HGA-1, and HGA-2 are all executed on that same instance using the same \(K\). The reported HGA configuration uses \texttt{popSize=100} and \texttt{generations=100}.

\subsection{Medium-Scale Setup}

The medium-scale bucket contains \(n\in\{50,100,150,200\}\) and the same ratio filter \(m/n\in\{2,4,6,8\}\). All available seed files are selected for each retained \((n,m)\) pair. As in the small-scale setting, \(K\) is first optimized separately for each instance using only \texttt{GRCCVC}. Then the final five-method comparison is carried out with the selected \(K\). In the reported configuration, the final HGA settings are \texttt{popSize=100} and \texttt{generations=100}. Importantly, the HGA methods are not used to choose \(K\); they are evaluated only after \(K\) has been fixed by the \texttt{GRCCVC}-based selection stage.

\subsection{Large-Scale Setup (Planned / Pending --- Placeholder)}

The large-scale bucket is defined in code as \(n\in\{250,500,750,1000\}\) with the same ratio filter \(m/n\in\{2,4,6,8\}\). If the same protocol is followed, \(K\) should first be selected per instance by the \texttt{GRCCVC}-based binary-search-plus-weight-scan stage and only then should the final five-method comparison be executed. However, by explicit instruction, no large-scale results are inserted in this draft.

\begin{table}[h]
\centering
\caption{Table A: Simulation Parameters}
\label{tab:sim-params}
\begin{tabular}{p{3.0cm}p{3.1cm}p{3.1cm}p{3.1cm}}
\toprule
Parameter & Small & Medium & Large (planned) \\
\midrule
Dataset source & DagdevirenDataset & DagdevirenDataset & DagdevirenDataset \\
Node counts & \{10,15,20,25\} & \{50,100,150,200\} & \{250,500,750,1000\} \\
Ratio filter & \{2,4,6,8\} & \{2,4,6,8\} & \{2,4,6,8\} \\
Selection mode & all available \(s\) per \((n,m)\) & all available \(s\) per \((n,m)\) & planned same policy \\
Methods & Greedy-1, Greedy-2, Greedy-3, HGA-1, HGA-2 & same & same \\
Exact solver & excluded by default & excluded by default & excluded by default \\
Capacity policy & optimized per instance via \texttt{GRCCVC} & optimized per instance via \texttt{GRCCVC} & planned same protocol \\
K search rule & binary search for min feasible \(K\) + sampled weight scan & same & planned same \\
K selector method & \texttt{GRCCVC} only & \texttt{GRCCVC} only & planned same \\
Final popSize & 100 & 100 & 100 planned \\
Final generations & 100 & 100 & 100 planned \\
Seed & common experiment seed & common experiment seed & planned same \\
HGA role in K search & not used & not used & planned not used \\
\bottomrule
\end{tabular}
\end{table}

\section{Results and Discussion}

\subsection{Small-Scale Results (Observed)}

Observed small-scale numerical results are not included here because they are maintained externally in the DOCX and were not provided as repository artifacts. To satisfy the integrity constraints, no numerical values are inserted below.

\begin{table}[h]
\centering
\caption{Table B: Method Comparison on Small-Scale Instances (Observed --- TO BE FILLED FROM DOCX)}
\label{tab:small-results}
\begin{tabular}{llllll}
\toprule
Method & Valid Runs & Valid Rate & Avg. Weight & Avg. Cover Size & Avg. Time (ms) \\
\midrule
Greedy-1 & TO BE FILLED & TO BE FILLED & TO BE FILLED & TO BE FILLED & TO BE FILLED \\
Greedy-2 & TO BE FILLED & TO BE FILLED & TO BE FILLED & TO BE FILLED & TO BE FILLED \\
Greedy-3 & TO BE FILLED & TO BE FILLED & TO BE FILLED & TO BE FILLED & TO BE FILLED \\
HGA-1 & TO BE FILLED & TO BE FILLED & TO BE FILLED & TO BE FILLED & TO BE FILLED \\
HGA-2 & TO BE FILLED & TO BE FILLED & TO BE FILLED & TO BE FILLED & TO BE FILLED \\
\bottomrule
\end{tabular}
\end{table}

The discussion should be completed only after the DOCX values are inserted. At that point, the comparison should address solution quality, cover cardinality, runtime trade-offs, and any feasibility failures.

\subsection{Medium-Scale Results (Observed)}

Observed medium-scale numerical results are likewise absent from the repository artifacts available in this session. Therefore, the following table is intentionally left empty.

\begin{table}[h]
\centering
\caption{Table C: Method Comparison on Medium-Scale Instances (Observed --- TO BE FILLED FROM DOCX)}
\label{tab:medium-results}
\begin{tabular}{llllll}
\toprule
Method & Valid Runs & Valid Rate & Avg. Weight & Avg. Cover Size & Avg. Time (ms) \\
\midrule
Greedy-1 & TO BE FILLED & TO BE FILLED & TO BE FILLED & TO BE FILLED & TO BE FILLED \\
Greedy-2 & TO BE FILLED & TO BE FILLED & TO BE FILLED & TO BE FILLED & TO BE FILLED \\
Greedy-3 & TO BE FILLED & TO BE FILLED & TO BE FILLED & TO BE FILLED & TO BE FILLED \\
HGA-1 & TO BE FILLED & TO BE FILLED & TO BE FILLED & TO BE FILLED & TO BE FILLED \\
HGA-2 & TO BE FILLED & TO BE FILLED & TO BE FILLED & TO BE FILLED & TO BE FILLED \\
\bottomrule
\end{tabular}
\end{table}

\subsection{Large-Scale Evaluation (Planned / Pending)}

No large-scale numerical values are inserted in this draft.

\begin{table}[h]
\centering
\caption{Table D: Large-Scale Results Template (TO BE FILLED)}
\label{tab:large-results}
\begin{tabular}{llllll}
\toprule
Method & Valid Runs & Valid Rate & Avg. Weight & Avg. Cover Size & Avg. Time (ms) \\
\midrule
Greedy-1 & TO BE FILLED & TO BE FILLED & TO BE FILLED & TO BE FILLED & TO BE FILLED \\
Greedy-2 & TO BE FILLED & TO BE FILLED & TO BE FILLED & TO BE FILLED & TO BE FILLED \\
Greedy-3 & TO BE FILLED & TO BE FILLED & TO BE FILLED & TO BE FILLED & TO BE FILLED \\
HGA-1 & TO BE FILLED & TO BE FILLED & TO BE FILLED & TO BE FILLED & TO BE FILLED \\
HGA-2 & TO BE FILLED & TO BE FILLED & TO BE FILLED & TO BE FILLED & TO BE FILLED \\
\bottomrule
\end{tabular}
\end{table}

\paragraph{Expected Scalability Trends (Hypotheses).}
The following points are expectations rather than observations.
\begin{itemize}
    \item Runtime is expected to increase with graph size for all methods.
    \item Greedy methods are expected to remain faster than HGA methods because they construct a single solution instead of evolving a population.
    \item HGA variants are expected to offer better weight-quality trade-offs on at least some denser instances because they search over decoder parameters.
    \item HGA-2 is expected to be more stable than HGA-1 on larger graphs due to its more controlled seeding, stochastic top-\(L\) decoder, and bounded local search.
    \item Capacity feasibility is expected to become more restrictive as density increases.
\end{itemize}

\section{Threats to Validity}

Several threats to validity should be acknowledged. First, the repository models secure monitoring at the graph level only. It does not explicitly implement IDS signatures, packet-content analysis, cryptographic trust establishment, radio interference, or attack-specific behavior. Consequently, secure monitoring in this work should be interpreted as connected, capacity-feasible link observability rather than a full-stack security evaluation.

Second, the semantic meaning of node weights is not explicitly documented in the provided artifacts. The code uses the weights exactly as stored in the dataset, but it does not specify whether they correspond to residual energy, monitoring cost, risk exposure, or another normalized quantity.

Third, external validity is limited by the graph abstraction. The workflow operates on weighted graph instances rather than on a physical simulator with channel effects, mobility, traffic classes, or industrial process semantics. Therefore, graph-level gains may not transfer quantitatively to all real deployments without further validation.

Fourth, internal validity depends on the experimental settings that were used when the external DOCX was produced. Although the repository exposes defaults and selection logic, any parameter overrides, execution hardware, and post-processing details are not specified in the provided artifacts.

\section{Missing Information for Full Reproducibility}

The following details are not specified in the provided artifacts and should be documented explicitly for full reproducibility.

\begin{itemize}
    \item The actual small-scale result values from the DOCX.
    \item The actual medium-scale result values from the DOCX.
    \item Confirmation that large-scale outputs are still pending.
    \item Whether the reported runs used repository defaults or parameter overrides for \texttt{capacityK}, \texttt{optimizeK}, \texttt{optimizeGoal}, \texttt{popSize}, \texttt{generations}, and \texttt{seed}.
    \item The exact command line or API payload used to generate the DOCX results.
    \item Hardware and software environment details for runtime measurements.
    \item Whether runtimes were collected from a single run or averaged across multiple runs.
    \item Any filtering, aggregation, or outlier-handling procedure applied before preparing the DOCX.
\end{itemize}

\clearpage
\selectlanguage{turkish}

\section*{Adım 1: Bulunan Gerçekleştirimler}

Depoda etkin olarak kullanılan beş yöntem aşağıdaki gibidir.

\begin{itemize}
    \item Açgözlü-1: \texttt{GCCVC}, giriş noktası \texttt{solve\_gccvc}, dosya \texttt{backend/app/algorithms.py}.
    \item Açgözlü-2: \texttt{GRCCVC}, giriş noktası \texttt{solve\_grccvc}, dosya \texttt{backend/app/algorithms.py}.
    \item Açgözlü-3: \texttt{GWCCVC}, giriş noktası \texttt{solve\_gwccvc}, dosya \texttt{backend/app/algorithms.py}.
    \item HGA-1: \texttt{HGA-CCVC}, giriş noktası \texttt{solve\_hga}, dosya \texttt{backend/app/algorithms.py}.
    \item HGA-2: \texttt{HGA-CCVC v2}, \texttt{HGA\_CCVC\_V2} sınıfı ve \texttt{solve\_hga\_v2} sarmalayıcısı, dosya \texttt{backend/app/algorithms.py}.
\end{itemize}

\section*{Adım 2: Koddan Çıkarılan Veri Kümesi Ayrımı ve Deney Protokolü}

Depo, \texttt{backend/data/DagdevirenDataset} altında tutulan tek bir veri kümesi kaynağını kullanmaktadır. Grafik dosyaları \texttt{n\{n\}\_m\{m\}\_s\{s\}.txt} biçimindedir. Koddaki varsayılan ölçek ayrımı
\[
\text{Küçük}=\{10,15,20,25\},\qquad
\text{Orta}=\{50,100,150,200\},\qquad
\text{Büyük}=\{250,500,750,1000\}
\]
şeklindedir. Varsayılan bağlantılılık oranı süzgeci
\[
\frac{m}{n}\in\{2,4,6,8\}
\]
olarak tanımlanmıştır. Ölçek bazlı varsayılan kapasite değerleri ise
\[
K_{\text{küçük}}=18,\qquad K_{\text{orta}}=16,\qquad K_{\text{büyük}}=16
\]
biçimindedir. Seçim politikası, seçilen her \((n,m)\) çifti için mevcut tüm \(s\) tohum dosyalarını içermektedir. Düzeltilmiş deney açıklamasına göre test süreci iki aşamada yürütülmüştür. İlk aşamada her grafik örneği için izleyici kapasitesi \(K\) otomatik olarak seçilmiştir. Arama aralığı
\[
K_{\min}=\left\lceil \frac{|E|}{|V|}\right\rceil,\qquad K_{\max}=\Delta(G)
\]
şeklinde belirlenmiştir. Bu aralıkta önce uygulanabilir en küçük \(K\) değerini bulmak için ikili arama yapılmış, bu aşamadaki uygulanabilirlik ve ağırlık seçicisi olarak oran tabanlı açgözlü yöntem \texttt{GRCCVC} kullanılmıştır. Amaç minimum ağırlık olduğunda, daha sonra uygulanabilir \(K\) değerleri üzerinde örneklenmiş bir tarama yapılmış ve geçerli \texttt{GRCCVC} ağırlığını en aza indiren \(K\) seçilmiştir; eşitlik durumunda daha küçük \(K\) tercih edilmiştir. İkinci aşamada ise bu örneğe özgü \(K\) sabitlendikten sonra aynı grafik üzerinde beş yöntemin tamamı çalıştırılmıştır. Raporlanan ayarlarda nüfus büyüklüğü ve nesil sayısı sırasıyla \texttt{100} ve \texttt{100} olarak kullanılmıştır.

\paragraph{Bütünlük notu.}
Bu dosyaya hiçbir sonuç değeri eklenmemiştir. Küçük ve orta ölçek sonuç tabloları, dışarıdaki DOCX dosyasından doldurulmak üzere özellikle boş bırakılmıştır. Büyük ölçek değerlendirmesi ise planlanan/bekleyen çalışma olarak işaretlenmiştir.

\section*{3. Sistem Modeli}

Endüstriyel kablosuz algılayıcı ağı (WSN), \(G=(V,E)\) biçiminde yönsüz bir grafik olarak modellenmektedir. Burada her \(v\in V\) düğümü bir sensör ya da izleme yapabilen düğümü, her \(e=(u,v)\in E\) kenarı ise bir haberleşme bağlantısını temsil eder. Gerçekleştirimde düğümler kimlik, ağırlık ve izleme kapasitesi taşırken; kenarlar kapsanmış/kapsanmamış durumu ve hangi düğüm tarafından izlendiği bilgisini taşır. Dolayısıyla izleme, yol tabanlı değil uç-düğüm tabanlıdır.

Bir \(e=(u,v)\) bağlantısı, uç noktalarından en az biri izleyici kümesine seçildiğinde ve söz konusu kenar seçilen bu uç noktaya atandığında izlenmiş kabul edilir. \(C\subseteq V\) seçilen izleyici kümesini göstersin. Bu durumda her \(e=(u,v)\in E\) için \(u\in C\) veya \(v\in C\) koşulu sağlanmalıdır; ayrıca her kapsanan kenar, seçilmiş komşu uçlardan tam olarak birine bağlanır.

Düğüm ağırlığı \(w(v)\), veri kümesinden doğrudan alınan dışsal bir maliyet değeridir. Kod, bu ağırlığın artık pil seviyesi, görev çevrimi ya da güvenlik maruziyetinden nasıl türetildiğini belirtmemektedir. Bu nedenle en doğru yorum, \(w(v)\) değerinin izleme yükünü veya enerji maliyetini temsil eden genel bir etkinleştirme maliyeti olduğudur. Amaç fonksiyonu, seçilen bağlı izleme kümesinin toplam ağırlığını en aza indirmektedir.

Her düğüme deney bazında aynı kapasite değeri \(K\) atanır:
\[
\mathrm{cap}(v)=K,\qquad \forall v\in V.
\]
Bu kapasite, \(v\) düğümüne atanabilecek komşu bağlantı sayısını sınırlar:
\[
\sum_{e\in \delta(v)} y_{e,v}\le K.
\]
Gerçekleştirimde her düğüm için artık kapasite tutulmakta ve düğüm bir kenarı üstlendiğinde bu değer azaltılmaktadır.

Bağlılık koşulu, izleyici düğümlerin ayrık kümeler yerine tek parça bir izleme omurgası oluşturmasını amaçlar. Bu nedenle geçerli bir çözümde \(C\) tarafından indüklenen altgrafın bağlı olması gerekir. Mevcut depoda ``güvenli izleme'' kavramı, bağlı ve kapasite açısından uygulanabilir bir izleyici kümesi ile tüm bağlantıların gözlemlenebilmesi şeklinde grafik düzeyinde temsil edilmektedir. Paket denetimi, IDS kuralları, kriptografik protokoller ve saldırgan modeli açık biçimde verilmemiştir.

\section*{4. Problem Formülasyonu (MWCCVC)}

\(x_v\in\{0,1\}\), \(v\in V\) düğümünün izleyici olarak seçilip seçilmediğini göstersin. Kapasite kısıtı altındaki kesin kenar-atama modelini ifade etmek için, yalnızca \(v\) düğümü \(e\) kenarının bir ucuysa tanımlı olan \(y_{e,v}\in\{0,1\}\) değişkenleri kullanılsın. Depodaki çözüm aşağıdaki amacı gerçekleştirmektedir:
\[
\min \sum_{v\in V} w(v)x_v.
\]

Kenar kapsama koşulu
\[
x_u+x_v\ge 1,\qquad \forall (u,v)\in E
\]
şeklindedir. Atama bakış açısıyla her kenar seçilmiş uçlarından tam birine atanmalıdır:
\[
\sum_{v\in e} y_{e,v}=1,\qquad \forall e\in E,
\]
ve tutarlılık kısıtı
\[
y_{e,v}\le x_v,\qquad \forall e\in E,\; \forall v\in e
\]
olarak yazılır.

Kapasite kısıtı ise
\[
\sum_{e\in \delta(v)} y_{e,v}\le Kx_v,\qquad \forall v\in V
\]
biçimindedir. Bu model koddaki gerçeklemeye yakındır; çünkü kapasite uygulanabilirliği, zorunlu atamalar ve esnek kenarlar için azami akış doğrulaması ile kontrol edilmektedir.

Bağlılık, \(C=\{v\in V:x_v=1\}\) kümesi için \(G[C]\) altgrafının bağlı olması şartıyla sağlanır:
\[
G[C]\text{ bağlıdır.}
\]
Kod, bunu tamsayılı kesme kısıtlarıyla değil; aday örtüyü oluşturduktan sonra bileşenleri en kısa köprü yolları ile onaran bir prosedürle sağlamaktadır.

Bu problem, ağırlıklı kapasiteli bağlı düğüm örtüsü türünde NP-zor bir optimizasyon problemidir [CITATION NEEDED]. Bu nedenle depo, açgözlü sezgiseller ve hibrit genetik algoritmalar kullanmaktadır.

\section*{5. Önerilen Yöntemler}

Beş yöntemin tamamı aynı uygulanabilirlik hedeflerini izler: tüm kenarların kapsanması, seçilen düğüm kümesinin bağlı olması ve kapasite kısıtının kesin olarak sağlanması. Üç açgözlü yöntem düğüm önceliklendirme ve kenar sıralama kuralları bakımından ayrışırken, iki HGA türü son çözümü doğrudan kodlamak yerine parametreli bir çözücü üzerinden arama yapar.

Gerçekleştirim ayrıca renk tabanlı bir yönlendirme mantığı kullanır: \texttt{WHITE} henüz işlenmemiş düğümü, \texttt{BLACK} seçilmiş izleyiciyi, \texttt{GRAY} seçilmiş izleyiciye komşu düğümü, \texttt{RED} ise öncelikli düğümü temsil eder.

\subsection*{5.1 Açgözlü-1}

Açgözlü-1, \texttt{GCCVC} yöntemine karşılık gelir. Yöntem derece odaklıdır: uygun adaylar arasından kapsanmamış derecesi en büyük olan düğüm seçilir; eşitlik halinde ağırlığı daha küçük olan tercih edilir. Bir düğüm seçildikten sonra, o düğümün kapsanmamış komşu kenarları, karşı uçtaki düğümün kapsanmamış derecesi azalan ve ardından komşu ağırlığı artan düzende işlenir. Böylece hem düğüm seçimi hem de yerel kenar ataması kapsama etkisini öncelemektedir.

Kapsama, seçilen düğümlere kenarların artık kapasite bitene kadar atanmasıyla artımlı biçimde sağlanır. Ana kurulum aşamasından sonra hâlâ kapsanmamış kenarlar varsa, yöntem önce seçilmiş düğümlerin kalan kapasitelerini tüketen, ardından gerekirse yeni düğümler ekleyen bir tamamlama aşaması kullanır.

Bağlılık, kapasite duyarlı bir onarım mekanizmasıyla sağlanır. Seçilen küme birden çok bileşen içeriyorsa, düşük maliyetli köprü yolları bulunur ve bileşenler birleştirilir.

\begin{algorithm}[h]
\caption{Açgözlü-1 / GCCVC}
\begin{algorithmic}[1]
\State \textbf{Girdi:} \(G=(V,E)\), ağırlıklar \(w(v)\), kapasite \(K\)
\State \textbf{Çıktı:} Bağlı kapasiteli örtü \(C\)
\State \(C\gets\emptyset\); tüm düğümleri \texttt{WHITE} olarak başlat; artık kapasiteleri \(K\) yap
\While{kapsanmamış bir kenar varsa}
    \State kapsanmamış derece kurallarıyla öncelikli (\texttt{RED}) düğümleri işaretle
    \State uygun aday kümesini oluştur
    \State kapsanmamış derecesi en büyük düğüm \(v\)'yi seç; eşitlikte ağırlığı küçük olanı al
    \State \(v\)'yi \(C\)'ye ekle ve \texttt{BLACK} yap
    \State \(v\)'nin kapsanmamış komşu kenarlarını uygun sıraya göre sırala
    \For{her sıralı kenar için}
        \If{\(v\)'nin artık kapasitesi sıfır ise}
            \State \textbf{break}
        \EndIf
        \State kenarı \(v\)'ye ata ve artık kapasiteyi azalt
    \EndFor
\EndWhile
\If{kapsanmamış kenar kaldıysa}
    \State kalan kenarları zorlayarak kapa
\EndIf
\State kopuk bileşenleri kapasite duyarlı köprü yolları ile birleştir
\State \Return \(C\)
\end{algorithmic}
\end{algorithm}

\subsection*{5.2 Açgözlü-2}

Açgözlü-2, \texttt{GRCCVC} yöntemidir. Düğüm seçme kuralı oran tabanlıdır:
\[
\rho(v)=\frac{w(v)}{\sum_{u\in N_U(v)} w(u)}.
\]
Yöntem, \(\rho(v)\) oranı en küçük olan düğümü seçer; eşitlik halinde daha düşük ağırlıklı düğüm tercih edilir. Böylece kendi etkinleştirme maliyeti, kapsayabildiği komşuluk yapısına göre düşük olan düğümler öne çıkar.

Seçilen düğümün kapsanmamış komşu kenarları atama öncesinde rastgele karıştırılır. Bu nedenle Açgözlü-2, düğüm seçiminde deterministik bir oran kuralını, yerel kenar atamasında ise çeşitlendirici bir rastgelelik ile birleştirir. Kapsama, tamamlama ve kapasite-duyarlı bağlılık onarımı Açgözlü-1 ile aynı genel mantığı izler.

\begin{algorithm}[h]
\caption{Açgözlü-2 / GRCCVC}
\begin{algorithmic}[1]
\State \textbf{Girdi:} \(G=(V,E)\), ağırlıklar \(w(v)\), kapasite \(K\)
\State \textbf{Çıktı:} Bağlı kapasiteli örtü \(C\)
\State \(C\gets\emptyset\) ve artık kapasiteleri \(K\) olarak başlat
\While{kapsanmamış bir kenar varsa}
    \State \texttt{RED} düğümleri işaretle
    \State aday kümesini oluştur
    \For{her aday \(v\) için}
        \State \(\rho(v)=w(v)/\sum_{u\in N_U(v)}w(u)\) hesapla
    \EndFor
    \State \(\rho(v)\) değeri en küçük olan düğümü seç; eşitlikte ağırlığı küçük olanı al
    \State \(v\)'yi \(C\)'ye ekle
    \State \(v\)'nin kapsanmamış komşu kenarlarını rastgele karıştır
    \State kenarları, \(v\)'nin artık kapasitesi bitene kadar sırayla \(v\)'ye ata
\EndWhile
\If{kapsanmamış kenar kaldıysa}
    \State kalan kenarları zorlayarak kapa
\EndIf
\State kopuk bileşenleri kapasite duyarlı köprü yolları ile onar
\State \Return \(C\)
\end{algorithmic}
\end{algorithm}

\subsection*{5.3 Açgözlü-3}

Açgözlü-3, \texttt{GWCCVC} yöntemidir. Bu yaklaşım ağırlık odaklıdır: minimum ağırlıklı aday düğüm seçilir; eşitlik durumunda kapsanmamış derecesi daha büyük olan tercih edilir. Böylece etkinleştirme maliyetine diğer açgözlü varyantlardan daha fazla ağırlık verilir.

Seçilen düğümün komşu kenarları sıralanırken, \texttt{WHITE} komşulara giden kenarlar önceliklendirilir. Bu tercih, daha önce hiç dokunulmamış bölgelere yayılmayı teşvik eder. Kapasite yine doğrudan kenar ataması sırasında uygulanır.

Kapsama tamamlaması ortak zorlama aşamasıyla yapılırken, bağlılık onarımı daha basit ve topolojiktir. Her bileşenden en küçük ağırlıklı bir temsilci seçilir ve bu temsilciler, bir temel bileşene BFS tabanlı en kısa yollarla bağlanır.

\begin{algorithm}[h]
\caption{Açgözlü-3 / GWCCVC}
\begin{algorithmic}[1]
\State \textbf{Girdi:} \(G=(V,E)\), ağırlıklar \(w(v)\), kapasite \(K\)
\State \textbf{Çıktı:} Bağlı kapasiteli örtü \(C\)
\State \(C\gets\emptyset\) ve artık kapasiteleri \(K\) olarak başlat
\While{kapsanmamış bir kenar varsa}
    \State \texttt{RED} düğümleri işaretle
    \State aday kümesini oluştur
    \State ağırlığı en küçük düğümü seç; eşitlikte kapsanmamış derecesi büyük olanı al
    \State düğümü \(C\)'ye ekle
    \State komşu kenarları, \texttt{WHITE} komşular öncelikli olacak biçimde sırala
    \State kenarları artık kapasite bitene kadar ata
\EndWhile
\If{kapsanmamış kenar kaldıysa}
    \State kalan kenarları zorlayarak kapa
\EndIf
\State \(G[C]\)'nin bileşenlerini bul
\State her bileşenden minimum ağırlıklı temsilci seç
\State temsilcileri BFS en kısa yolları ile bir temel bileşene bağla
\State köprü yolundaki düğümleri \(C\)'ye ekle
\State \Return \(C\)
\end{algorithmic}
\end{algorithm}

\subsection*{5.4 HGA-1}

HGA-1, \texttt{HGA-CCVC} yöntemine karşılık gelir. Kromozom iki bölümden oluşur: altı boyutlu gerçek değerli parametre vektörü
\[
\boldsymbol{\theta}=(a,b,c,d,e,\epsilon)\in[0,3]^6
\]
ve düğüm bazlı rastgele anahtar değerleri. Kromozom doğrudan son örtüyü kodlamaz; bunun yerine baştan örtü oluşturan parametreli bir açgözlü çözücüyü yönlendirir.

Bir aday düğüm için çözücü skoru şu şekilde özetlenebilir:
\[
S(v)=a\deg_U(v)-bw(v)-c\rho(v)+d\mathbf{1}_{\mathrm{RED}}(v)+e\mathbf{1}_{\mathrm{GRAY}}(v)+0.3a\frac{\deg_U(v)}{\max(w(v),10^{-9})}-\epsilon\,\mathrm{key}(v).
\]
Burada \(\deg_U(v)\) kapsanmamış dereceyi, \(\rho(v)\) ise oran terimini göstermektedir. Böylece yöntem, kapsama etkisi, düğüm maliyeti, komşuluk yapısı, renk bonusları ve çeşitlendirici rastgeleliği birlikte dengeleyen parametreleri aramaktadır.

Başlatma hibrit yapıdadır. İlk nüfus, üç açgözlü yöntemin ürettiği örtüler ile ön tanımlı \(\theta\) şablonlarını birleştirerek tohumlanır. Geri kalan bireyler, bu şablonların pertürbe edilmesi veya tamamen rastgele oluşturulmasıyla elde edilir.

Kısıt işleme de hibrittir. Çözücü önce bir örtü inşa eder ve bağlılığı onarır; ardından geçersiz çözümler cezalandırılır. Ceza fonksiyonu özetle
\[
f(C)=W(C)+10^6U(C)+10^5(1-I_{\mathrm{conn}}(C))+10^5(1-I_{\mathrm{cap}}(C))
\]
şeklindedir. Burada \(W(C)\) toplam ağırlığı, \(U(C)\) kapsanmamış kenar sayısını göstermektedir.

HGA-1 ayrıca budama ve ağırlık odaklı yerel arama, elitizm, turnuva seçimi, çaprazlama, mutasyon ve durgunluk halinde yeni tohum birey enjeksiyonu içermektedir.

\begin{algorithm}[h]
\caption{HGA-1 / HGA-CCVC}
\begin{algorithmic}[1]
\State \textbf{Girdi:} \(G=(V,E)\), ağırlıklar, kapasite \(K\), nüfus büyüklüğü \(P\), nesil sayısı \(G_{\max}\)
\State \textbf{Çıktı:} En iyi uygulanabilir bağlı kapasiteli örtü \(C^{\star}\)
\State hibrit başlangıç nüfusunu oluştur
\State \(C^{\star}\gets\emptyset\), \(f^{\star}\gets +\infty\)
\For{her nesil için}
    \For{her kromozom için}
        \State kromozomu parametreli açgözlü çözücü ile bir örtüye çöz
        \State bağlılığı onar
        \State isteğe bağlı yerel arama uygula
        \State cezalı uygunluk değerini hesapla
        \If{daha iyi ise}
            \State en iyi çözümü güncelle
        \EndIf
    \EndFor
    \State nüfusu sırala ve elitleri koru
    \If{durgunluk oluştuysa}
        \State yeni tohum birey ekle
    \EndIf
    \While{yeni nüfus dolmadıysa}
        \State ebeveynleri turnuva ile seç
        \State çaprazlama ve mutasyon uygula
        \State yavruyu nüfusa ekle
    \EndWhile
\EndFor
\State \Return \(C^{\star}\)
\end{algorithmic}
\end{algorithm}

\subsection*{5.5 HGA-2}

HGA-2, \texttt{HGA-CCVC v2} yöntemidir. Temel kromozom fikri HGA-1 ile aynıdır; ancak evrimsel tasarım daha modüler ve daha açık parametrelerle gerçekleştirilmiştir. Kodda varsayılan olarak \texttt{elite\_size=2}, \texttt{tournament\_k=3}, \texttt{theta\_mut\_prob=0.2}, \texttt{theta\_mut\_sigma=0.2}, \texttt{key\_mut\_prob=0.02}, \texttt{local\_search\_prob=0.2}, \texttt{local\_search\_moves=7}, \texttt{local\_search\_trials=36}, \texttt{top\_l=3}, \texttt{repair\_mode=tier2}, \texttt{stall\_generations=10} ve \texttt{seed\_bias=0.3} kullanılmaktadır.

Başlatma hibrit ve tohumludur. Bir \texttt{GCCVC}, bir \texttt{GRCCVC} ve iki \texttt{GWCCVC} tabanlı strateji şablonu kullanılır; ardından hedef nüfus büyüklüğüne kadar rastgele bireyler eklenir.

Çözücü skor tabanlıdır; ancak HGA-2, üst-\(L\) stokastik seçim mekanizması ekler. Varsayılan durumda her zaman en yüksek skorlu düğüm seçilmez; en iyi birkaç aday arasından örnekleme yapılır. Benzer bir çeşitlendirme yerel kenar sıralamasına da uygulanır.

Ceza fonksiyonu özetle
\[
f(C)=W(C)+0.001|C|+\lambda_1U(C)+\lambda_2(1-I_{\mathrm{conn}}(C))+\lambda_3\phi_{\mathrm{cap}}(C)
\]
şeklindedir. Burada \(\lambda_1=10^6\), \(\lambda_2=10^5\) ve \(\lambda_3=10^5\) alınır. \(0.001|C|\) terimi, benzer ağırlıklı çözümler arasında daha küçük örtüyü tercih ettiren bağlayıcı bir terimdir.

Seçim \(k\)-turnuva ile, \(\theta\) üzerinde aritmetik çaprazlama ve anahtarlar üzerinde uniform çaprazlama ile yapılır. Mutasyon ise \(\theta\) için Gauss gürültüsü, anahtarlar için rastgele yenileme içerir. Yerel arama yalnızca uygulanabilir çözümleri kabul eder.

\begin{algorithm}[h]
\caption{HGA-2 / HGA-CCVC v2}
\begin{algorithmic}[1]
\State \textbf{Girdi:} \(G=(V,E)\), ağırlıklar, kapasite \(K\), nüfus büyüklüğü \(P\), nesil sayısı \(G_{\max}\)
\State \textbf{Çıktı:} En iyi uygulanabilir bağlı kapasiteli örtü \(C^{\star}\)
\State nüfusu açgözlü tabanlı tohumlar ve rastgele bireylerle başlat
\For{her nesil için}
    \For{her kromozom için}
        \State skor güdümlü açgözlü çözücü ile örtüyü oluştur
        \State stokastik kipte üst-\(L\) adaydan örnekle
        \State bağlılığı seçilen onarım kipine göre düzelt
        \State uygulanabilir-odaklı yerel arama uygula
        \State cezalı uygunluğu hesapla
    \EndFor
    \State bireyleri sırala ve elitleri koru
    \If{durgunluk eşiğine ulaşıldıysa}
        \State çalışmayı sonlandır
    \EndIf
    \While{nüfus dolmadıysa}
        \State ebeveynleri turnuva ile seç
        \State çaprazlama ve mutasyon uygula
        \State yavruları ekle
    \EndWhile
\EndFor
\State \Return en iyi uygulanabilir çözüm
\end{algorithmic}
\end{algorithm}

\section*{6. Algoritmik Karmaşıklık Analizi}

\(n=|V|\), \(m=|E|\), \(P\) nüfus büyüklüğü, \(G\) nesil sayısı olsun. Kodun kullandığı kesin kapasite doğrulaması, esnek kenarlar ve seçili düğümler üzerinde bir akış ağı kurmaktadır. Bu maliyet \(T_{\mathrm{flow}}(n,m)\) ile gösterilsin.

Üç açgözlü yöntemde baskın maliyet; aday düğümlerin tekrar tekrar puanlanması, komşu kenarların taranması ve bağlılık onarımıdır. Yapıcı aşama için temkinli üst sınır \(O(nm)\) olarak alınabilir.

Açgözlü-1 ve Açgözlü-2 için kapasite duyarlı bağlılık onarımı nedeniyle temkinli sınır
\[
T_{\text{Açgözlü-1}}=O\bigl(nm+\kappa^2m\log n\bigr),
\]
\[
T_{\text{Açgözlü-2}}=O\bigl(nm+\kappa^2m\log n\bigr)
\]
şeklindedir. Burada \(\kappa\), onarım öncesi bileşen sayısını göstermektedir.

Açgözlü-3 için BFS tabanlı daha basit onarım nedeniyle sınır
\[
T_{\text{Açgözlü-3}}=O\bigl(nm+\kappa(n+m)\bigr)
\]
olarak verilebilir.

HGA-1 için her kromozom değerlendirmesi; çözücü çalıştırma, bağlılık onarımı, yerel arama ve kapasite doğrulaması içerir. Yerel arama maliyeti \(L_1\) ile gösterilirse,
\[
T_{\text{HGA-1}}=O\Bigl(PG\bigl(nm+\kappa^2m\log n+L_1T_{\mathrm{flow}}(n,m)\bigr)\Bigr)
\]
elde edilir.

HGA-2 için benzer biçimde, sınırlandırılmış yerel arama maliyeti \(L_2\) ile,
\[
T_{\text{HGA-2}}=O\Bigl(PG\bigl(nm+\kappa^2m\log n+L_2T_{\mathrm{flow}}(n,m)\bigr)\Bigr)
\]
şeklinde yazılabilir. Bu sınırlar uygulamadaki mevcut veri yapıları için temkinli karmaşıklık ifadeleridir.

\section*{7. Veri Kümesi Tanımı (Küçük/Orta/Büyük ayrımı)}

Çalışma, \texttt{backend/data/DagdevirenDataset} altında bulunan tek bir topoloji kaynağını kullanmaktadır. Dosyalar \((n,m,s)\) ile tanımlanır; burada \(n\) düğüm sayısını, \(m\) kenar sayısını, \(s\) ise tohum/örnek indeksini ifade eder. Okuyucu, isteğe bağlı \texttt{n m} başlık satırını destekler; başlık yoksa değerler dosya adından çıkarılır.

Başlıktan sonra gövde yapısı şöyledir:
\begin{enumerate}
    \item \texttt{id weight} biçiminde \(n\) adet düğüm kaydı,
    \item \texttt{u v} biçiminde \(m\) adet kenar kaydı.
\end{enumerate}

Varsayılan ölçek ayrımı
\[
\text{Küçük}=\{10,15,20,25\},\qquad
\text{Orta}=\{50,100,150,200\},\qquad
\text{Büyük}=\{250,500,750,1000\}
\]
şeklindedir. Varsayılan bağlantılılık oranı filtresi ise
\[
\frac{m}{n}\in\{2,4,6,8\}
\]
olarak tanımlıdır. Ön-ayarlı seçim politikası, seçilen her \((n,m)\) çifti için mevcut tüm tohum dosyalarını içermektedir.

\begin{table}[h]
\centering
\caption{Veri Kümesi Biçimi}
\label{tab:dataset-format-tr}
\begin{tabular}{lll}
\toprule
Alan & Anlamı & Örnek \\
\midrule
\texttt{n100\_m200\_s1.txt} & Dosya adı \((n,m,s)\) bilgisini kodlar & \texttt{n100\_m200\_s1.txt} \\
\texttt{n m} & İsteğe bağlı başlık satırı & \texttt{100 200} \\
\texttt{id weight} & Düğüm kaydı & \texttt{0 116.20978281782678} \\
\texttt{u v} & Kenar kaydı & \texttt{0 1} \\
\bottomrule
\end{tabular}
\end{table}

\section*{8. Deneysel Kurulum ve Protokol}

Tüm yöntemler, aynı ölçek kümesindeki aynı grafik örnekleri üzerinde, aynı düğüm ağırlıkları, aynı kenar kümesi ve o örnek için seçilmiş aynı kapasite değeri ile çalıştırılır. Uygulanabilirlik tek bir doğrulama yordamı ile denetlenir; bu yordam kenar kapsaması, bağlılık, toplam seçili ağırlık, örtü büyüklüğü ve kesin uç-atamalı kapasite uygulanabilirliğini raporlar.

Beş yöntemli değerlendirme hattı
\[
\{\texttt{gccvc},\texttt{grccvc},\texttt{gwccvc},\texttt{hga},\texttt{hga\_v2}\}
\]
şeklindedir. Düzeltilmiş protokolde \(K\) ölçek bazında önceden sabitlenmez. Bunun yerine her grafik örneği için
\[
K_{\min}=\left\lceil \frac{|E|}{|V|}\right\rceil,\qquad K_{\max}=\Delta(G)
\]
aralığı oluşturulur. Ardından \(K\)-seçim aşamasında yalnızca \texttt{GRCCVC} kullanılır. Önce uygulanabilir en küçük \(K\) değeri ikili arama ile bulunur. Daha sonra minimum ağırlık hedefi altında, uygulanabilir en küçük \(K\) ile \(K_{\max}\) arasındaki örneklenmiş \(K\) değerleri taranır. Seçilen kapasite, geçerli \texttt{GRCCVC} ağırlığını en aza indiren değerdir; eşitlikte daha küçük \(K\) alınır. Bu örneğe özgü \(K\) belirlendikten sonra beş yöntemin tamamı aynı grafik üzerinde bu \(K\) ile çalıştırılır. Raporlanan deneylerde \texttt{popSize=100}, \texttt{generations=100} ve ortak bir tohum ayarı kullanılmıştır.

\subsection*{8.1 Küçük Ölçek Kurulumu}

Küçük ölçek kümesi \(n\in\{10,15,20,25\}\) ve \(m/n\in\{2,4,6,8\}\) koşullarını sağlar. Seçilen her \((n,m)\) çifti için mevcut tüm tohum dosyaları kullanılır. Her küçük ölçek örneği için önce yukarıda açıklanan \texttt{GRCCVC} güdümlü ikili-arama artı ağırlık-tarama prosedürü ile \(K\) optimize edilir. Daha sonra Greedy-1, Greedy-2, Greedy-3, HGA-1 ve HGA-2 aynı örnek üzerinde seçilmiş bu \(K\) ile çalıştırılır. Raporlanan HGA ayarları \texttt{popSize=100} ve \texttt{generations=100} değerleridir.

\subsection*{8.2 Orta Ölçek Kurulumu}

Orta ölçek kümesi \(n\in\{50,100,150,200\}\) ve yine \(m/n\in\{2,4,6,8\}\) filtrelerini içerir. Seçilen her \((n,m)\) çifti için tüm mevcut tohum dosyaları kullanılır. Küçük ölçekte olduğu gibi, \(K\) her örnek için önce yalnızca \texttt{GRCCVC} kullanılarak optimize edilir. Sonrasında seçilmiş \(K\) ile beş yöntemin nihai karşılaştırması yapılır. Raporlanan ayarlarda son HGA çalıştırmaları için \texttt{popSize=100} ve \texttt{generations=100} kullanılmıştır. Önemli nokta şudur: HGA yöntemleri \(K\) seçiminde kullanılmaz; yalnızca \texttt{GRCCVC}-tabanlı seçim aşaması tamamlandıktan sonra değerlendirilirler.

\subsection*{8.3 Büyük Ölçek Kurulumu (Planlanan / Bekleyen --- Yer Tutucu)}

Büyük ölçek kümesi, kodda \(n\in\{250,500,750,1000\}\) ve oran filtresi \(m/n\in\{2,4,6,8\}\) olacak şekilde tanımlanmıştır. Aynı protokol izlenecekse, önce her örnek için \texttt{GRCCVC}-tabanlı ikili arama artı ağırlık taraması ile \(K\) seçilmeli, daha sonra beş yöntemin nihai karşılaştırması yapılmalıdır. Ancak açık talimat gereği bu taslağa hiçbir büyük ölçek sonuç değeri eklenmemiştir.

\begin{table}[h]
\centering
\caption{Tablo A: Simülasyon Parametreleri}
\label{tab:sim-params-tr}
\begin{tabular}{p{3.0cm}p{3.1cm}p{3.1cm}p{3.1cm}}
\toprule
Parametre & Küçük & Orta & Büyük (planlanan) \\
\midrule
Veri kaynağı & DagdevirenDataset & DagdevirenDataset & DagdevirenDataset \\
Düğüm sayıları & \{10,15,20,25\} & \{50,100,150,200\} & \{250,500,750,1000\} \\
Oran filtresi & \{2,4,6,8\} & \{2,4,6,8\} & \{2,4,6,8\} \\
Seçim kipi & her \((n,m)\) için tüm \(s\) & her \((n,m)\) için tüm \(s\) & planlanan aynı politika \\
Yöntemler & Açgözlü-1, Açgözlü-2, Açgözlü-3, HGA-1, HGA-2 & aynı & aynı \\
Kesin çözücü & varsayılan olarak hariç & varsayılan olarak hariç & varsayılan olarak hariç \\
Kapasite politikası & örnek bazında \texttt{GRCCVC} ile optimize edilir & örnek bazında \texttt{GRCCVC} ile optimize edilir & planlanan aynı protokol \\
K arama kuralı & uygulanabilir en küçük \(K\) için ikili arama + örneklenmiş ağırlık taraması & aynı & planlanan aynı \\
K seçici yöntem & yalnızca \texttt{GRCCVC} & yalnızca \texttt{GRCCVC} & planlanan aynı \\
Son popSize & 100 & 100 & planlanan 100 \\
Son generations & 100 & 100 & planlanan 100 \\
Tohum & ortak deney tohumu & ortak deney tohumu & planlanan aynı \\
HGA'nın K aramadaki rolü & kullanılmaz & kullanılmaz & planlanan kullanılmaz \\
\bottomrule
\end{tabular}
\end{table}

\section*{9. Sonuçlar ve Tartışma}

\subsection*{9.1 Küçük Ölçek Sonuçları (Gözlenen)}

Küçük ölçek sayısal sonuçlar bu oturumda sağlanan depo çıktıları içinde yer almamaktadır; dış DOCX dosyasından gelecektir. Bütünlük kuralı gereği aşağıya hiçbir sayısal değer yazılmamıştır.

\begin{table}[h]
\centering
\caption{Tablo B: Küçük Ölçek Yöntem Karşılaştırması (Gözlenen --- DOCX'TEN DOLDURULACAK)}
\label{tab:small-results-tr}
\begin{tabular}{llllll}
\toprule
Yöntem & Geçerli Çalışma & Geçerlilik Oranı & Ort. Ağırlık & Ort. Örtü Boyutu & Ort. Süre (ms) \\
\midrule
Açgözlü-1 & TO BE FILLED & TO BE FILLED & TO BE FILLED & TO BE FILLED & TO BE FILLED \\
Açgözlü-2 & TO BE FILLED & TO BE FILLED & TO BE FILLED & TO BE FILLED & TO BE FILLED \\
Açgözlü-3 & TO BE FILLED & TO BE FILLED & TO BE FILLED & TO BE FILLED & TO BE FILLED \\
HGA-1 & TO BE FILLED & TO BE FILLED & TO BE FILLED & TO BE FILLED & TO BE FILLED \\
HGA-2 & TO BE FILLED & TO BE FILLED & TO BE FILLED & TO BE FILLED & TO BE FILLED \\
\bottomrule
\end{tabular}
\end{table}

Bu alt bölüm, DOCX değerleri eklendikten sonra çözüm kalitesi, örtü büyüklüğü, çalışma süresi ve uygulanabilirlik başarısı bakımından tamamlanmalıdır.

\subsection*{9.2 Orta Ölçek Sonuçları (Gözlenen)}

Orta ölçek sayısal sonuçlar da bu oturumda sağlanan depo çıktıları içinde bulunmadığından tablo bilinçli olarak boş bırakılmıştır.

\begin{table}[h]
\centering
\caption{Tablo C: Orta Ölçek Yöntem Karşılaştırması (Gözlenen --- DOCX'TEN DOLDURULACAK)}
\label{tab:medium-results-tr}
\begin{tabular}{llllll}
\toprule
Yöntem & Geçerli Çalışma & Geçerlilik Oranı & Ort. Ağırlık & Ort. Örtü Boyutu & Ort. Süre (ms) \\
\midrule
Açgözlü-1 & TO BE FILLED & TO BE FILLED & TO BE FILLED & TO BE FILLED & TO BE FILLED \\
Açgözlü-2 & TO BE FILLED & TO BE FILLED & TO BE FILLED & TO BE FILLED & TO BE FILLED \\
Açgözlü-3 & TO BE FILLED & TO BE FILLED & TO BE FILLED & TO BE FILLED & TO BE FILLED \\
HGA-1 & TO BE FILLED & TO BE FILLED & TO BE FILLED & TO BE FILLED & TO BE FILLED \\
HGA-2 & TO BE FILLED & TO BE FILLED & TO BE FILLED & TO BE FILLED & TO BE FILLED \\
\bottomrule
\end{tabular}
\end{table}

\subsection*{9.3 Büyük Ölçek Değerlendirmesi (Planlanan / Bekleyen)}

Bu taslakta büyük ölçek için hiçbir sayısal sonuç verilmemiştir.

\begin{table}[h]
\centering
\caption{Tablo D: Büyük Ölçek Sonuç Şablonu (TO BE FILLED)}
\label{tab:large-results-tr}
\begin{tabular}{llllll}
\toprule
Yöntem & Geçerli Çalışma & Geçerlilik Oranı & Ort. Ağırlık & Ort. Örtü Boyutu & Ort. Süre (ms) \\
\midrule
Açgözlü-1 & TO BE FILLED & TO BE FILLED & TO BE FILLED & TO BE FILLED & TO BE FILLED \\
Açgözlü-2 & TO BE FILLED & TO BE FILLED & TO BE FILLED & TO BE FILLED & TO BE FILLED \\
Açgözlü-3 & TO BE FILLED & TO BE FILLED & TO BE FILLED & TO BE FILLED & TO BE FILLED \\
HGA-1 & TO BE FILLED & TO BE FILLED & TO BE FILLED & TO BE FILLED & TO BE FILLED \\
HGA-2 & TO BE FILLED & TO BE FILLED & TO BE FILLED & TO BE FILLED & TO BE FILLED \\
\bottomrule
\end{tabular}
\end{table}

\paragraph{Beklenen Ölçeklenebilirlik Eğilimleri (Hipotezler).}
Aşağıdaki maddeler gözlem değil, beklentidir.
\begin{itemize}
    \item Grafik boyutu arttıkça tüm yöntemlerin çalışma süresinin artması beklenir.
    \item Açgözlü yöntemlerin, nüfus tabanlı evrimsel arama yapmadıkları için HGA yöntemlerinden daha hızlı kalması beklenir.
    \item HGA türlerinin, özellikle daha yoğun örneklerde, ağırlık kalitesi bakımından daha iyi ödünleşimler sunması beklenir.
    \item HGA-2'nin, daha kontrollü tohumlama ve sınırlandırılmış yerel arama nedeniyle büyük grafiklerde HGA-1'e göre daha kararlı davranması beklenir.
    \item Yoğunluk arttıkça kapasite uygulanabilirliğinin daha belirleyici hâle gelmesi beklenir.
\end{itemize}

\section*{10. Geçerlilik Tehditleri}

Öncelikle, depo güvenli izlemeyi yalnızca grafik düzeyinde modellemektedir. IDS imzaları, paket içeriği analizi, kriptografik güven tesis mekanizmaları, girişim etkileri veya saldırı türlerine özgü davranışlar açık biçimde işlenmemektedir. Dolayısıyla bu çalışmadaki güvenli izleme, tam yığın bir güvenlik değerlendirmesi değil; bağlı ve kapasite açısından uygulanabilir bağlantı gözlenebilirliği olarak anlaşılmalıdır.

İkinci olarak, düğüm ağırlıklarının anlamsal yorumu açık biçimde belgelenmemiştir. Kod, ağırlıkları veri kümesinde verildiği şekliyle kullanmaktadır; ancak bunların artık enerji, izleme maliyeti, risk maruziyeti ya da başka bir normalize büyüklük olup olmadığı belirtilmemektedir.

Üçüncü olarak, dış geçerlilik grafik soyutlamasıyla sınırlıdır. İş akışı, kanal etkileri, hareketlilik, trafik türleri veya endüstriyel süreç semantiği içeren fiziksel bir benzetim yerine ağırlıklı grafik örnekleri üzerinde çalışmaktadır.

Son olarak, iç geçerlilik, dış DOCX hazırlanırken kullanılan deney ayarlarına bağlıdır. Kod varsayılanları ve seçim mantığı bilinse de, parametre üstten yazımları, donanım ortamı ve sonradan yapılan toplulaştırma ayrıntıları sağlanan artifaktlarda mevcut değildir.

\section*{11. Tam Yeniden Üretilebilirlik İçin Eksik Bilgiler}

Tam yeniden üretilebilirlik için aşağıdaki bilgiler sağlanan artifaktlarda belirtilmemiştir ve açıkça belgelenmelidir.

\begin{itemize}
    \item DOCX dosyasındaki gerçek küçük ölçek sonuç değerleri.
    \item DOCX dosyasındaki gerçek orta ölçek sonuç değerleri.
    \item Büyük ölçek çıktılarının gerçekten beklemede olduğuna dair teyit.
    \item Raporlanan deneylerin \texttt{capacityK}, \texttt{optimizeK}, \texttt{optimizeGoal}, \texttt{popSize}, \texttt{generations} ve \texttt{seed} için varsayılan ayarlarla mı yoksa üstten yazılmış parametrelerle mi çalıştırıldığı.
    \item DOCX sonuçlarını üretmek için kullanılan tam komut satırı veya API yükü.
    \item Süre ölçümleri için donanım ve yazılım ortamı bilgileri.
    \item Çalışma sürelerinin tek koşudan mı yoksa çoklu koşuların ortalamasından mı üretildiği.
    \item DOCX hazırlanırken uygulanan filtreleme, toplulaştırma veya aykırı değer işleme prosedürü.
\end{itemize}

\end{document}
